\section{Introduction}
\label{sec:introduction}

A neural vocoder reconstructs speech from an acoustic representation.
Autoregressive systems generate samples sequentially
\citep{oord2016wavenet}, whereas adversarial, inverse short-time Fourier
transform (iSTFT), and flow models move learned computation toward parallel
waveform synthesis, frame-level spectral prediction, or a few solver steps
\citep{kong2020hifigan,siuzdak2024vocos,liu2025rfwave}.  These computational
graphs expose different bottlenecks, so parameter count, serial dependency,
memory traffic, transform overhead, and solver evaluations need not track
one another.

Published quality-efficiency rankings also mix corpora, preprocessing,
training budgets, checkpoint rules, implementations, evaluators, and
hardware.  Even within one implementation, parameter count, serialized
bytes, multiply-accumulate counts, and real-time factor (RTF) answer
different questions.  Deployment transformations widen the gap between a
method name and the artifact that executes: compilation can break a graph,
quantization can cover only selected operators, and dense pruning can retain
neither sparse storage nor sparse kernels
\citep{han2015pruning,jacob2018quantization,ansel2024pytorch2}.

This work benchmarks twelve neural-vocoder configurations trained from
scratch under identical data splits and one evaluation implementation.  The
cohort comprises APNet2, BigVGAN-base, FreeV, HiFi-GAN V1, V2, V3, LPCNet,
MelGAN, RFWave, RNDVoC, Vocos, and VocosFormer.  Architecture-native
conditioning, objectives, and training budgets are retained.  Each configuration
contributes one training trajectory from recorded seed 1234 and one retained
checkpoint.  Inference seeds 42, 43, and 44 re-execute that frozen state
rather than retrain; the single exception is static ONNX INT8, which
rebuilds calibration state on each execution.

Three research questions structure the study.  \emph{RQ1}: under identical
data splits and one evaluation implementation, but
architecture-native objectives and budgets, how do the twelve
configurations trade objective quality, inference latency,
parameter count, and deployable storage?  \emph{RQ2}: relative to
same-hardware baselines, how do executed deployment transformations shift
objective quality, warm real-time factor, and serialized artifact size,
and how strongly do those effects depend on architecture and execution
backend?  \emph{RQ3}: which measured deployment artifact is selected under
explicit hardware-specific constraints on quality, parameter count,
and deployable storage?  The three are dependent analyses of one
experiment: transformed variants inherit the trained checkpoints, selection
consumes their measurements, and unsupported transformations are excluded
on record.

The scope is deliberately conditional.  LJSpeech contains one speaker.  The
final 525 utterances form an adaptive evaluation set: their PESQ values
informed checkpoint decisions and one deterministic repair of RFWave's
inference sampler.
Objective measures are proxies, not listener judgments
\citep{rix2001pesq,taal2011stoi,saeki2022utmos}.  Population-level or causal claims are precluded by one training trajectory
per configuration, unequal architecture-native budgets, independently
scheduled training jobs, and the absence of listening or multi-speaker
experiments.
